\subsection{FREE ALGEBRAS}

DEFINITION 2. Let I be a set; let M(I) (resp. Mo(I), resp. N \(^{(I)}\) ) denote the free magma (resp. free monoid, resp. free commutative monoid) derived from I. The algebra of M(I) (resp. Mo(I), resp. N \(^{(I)}\) ) over A is called the free algebra (resp. free associative algebra, resp. free commutative associative algebra (or, by an abuse of language, free commutative algebra)) of the set I over the ring A.

We shall denote the free algebra (resp. free associative algebra, resp. free commutative algebra) of I over A by \(\operatorname{Lib}_{\mathbf{A}}(\mathbf{I})\) (resp. \(\operatorname{Libas}_{\mathbf{A}}(\mathbf{I})\) , resp. \(\operatorname{Libasc}_{\mathbf{A}}(\mathbf{I})\) ). By composing the canonical mapping of I into \(\operatorname{M}(\mathbf{I})\) (resp. \(\operatorname{Mo}(\mathbf{I})\) , resp. \(\mathbf{N}^{(\mathbf{I})}\) ) with the canonical mapping of \(\operatorname{M}(\mathbf{I})\) (resp. \(\operatorname{Mo}(\mathbf{I})\) , resp. \(\mathbf{N}^{(\mathbf{I})}\) ) into \(\operatorname{Lib}_{\mathbf{A}}(\mathbf{I})\) (resp. \(\operatorname{Libas}_{\mathbf{A}}(\mathbf{I})\) , resp. \(\operatorname{Libasc}_{\mathbf{A}}(\mathbf{I})\) ), a canonical mapping is obtained of I into \(\operatorname{Lib}_{\mathbf{A}}(\mathbf{I})\) (resp. \(\operatorname{Libas}_{\mathbf{A}}(\mathbf{I})\) , resp. \(\operatorname{Libasc}_{\mathbf{A}}(\mathbf{I})\) ), which is injective if \(A \neq \{0\}\) . We shall denote the image of an element \(i \in I\) under this canonical mapping by \(X_{i}\) and we shall say that \(X_{i}\) is the indeterminate of index i of \(\operatorname{Lib}_{\mathbf{A}}(\mathbf{I})\) (resp. \(\operatorname{Libas}_{\mathbf{A}}(\mathbf{I})\) , resp. \(\operatorname{Libasc}_{\mathbf{A}}(\mathbf{I})\) ).

As \(\mathbf{Mo}(\mathbf{I})\) and \(\mathbf{N}^{(\mathrm{I})}\) each have an identity element, \(\operatorname{Libas}_{\mathbf{A}}(\mathbf{I})\) and \(\operatorname{Libasc}_{\mathbf{A}}(\mathbf{I})\)

are unital associative algebras and further \(\operatorname{Libasc}_{\mathbf{A}}(\mathbf{I})\) is commutative. If e is the unit element of \(\operatorname{Libas}_{\mathbf{A}}(\mathbf{I})\) (resp. \(\operatorname{Libasc}_{\mathbf{A}}(\mathbf{I})\) ), the mapping \(\alpha \mapsto \alpha e\) is an isomorphism of A onto a subring of the centre of \(\operatorname{Libas}_{\mathbf{A}}(\mathbf{I})\) (resp. \(\operatorname{Libasc}_{\mathbf{A}}(\mathbf{I})\) ), which is identified with A (no. 1).

PROPOSITION 7. Let I be a set and F an algebra (resp. unital associative algebra, resp. unital commutative associative algebra) over A. For every mapping \(f:\mathbf{I}\to\mathbf{F}\) , there exists one and only one homomorphism (resp. unital homomorphism) \(\bar{f}\) of \(\operatorname{Lib}_{\mathbf{A}}(\mathbf{I})\) (resp. \(\operatorname{Libas}_{\mathbf{A}}(\mathbf{I})\) , resp. \(\operatorname{Libasc}_{\mathbf{A}}(\mathbf{I})\) ) into F such that \(\bar{f}(\mathbf{X}_{i})=f(i)\) for all \(i\in\mathbf{I}\) .

Let \(F_{m}\) be the magma (resp. monoid) obtained by giving the set F its multiplicative law of composition. There is one and only one homomorphism (resp. unital homomorphism) g of M(I) (resp. Mo(I), resp. N \(^{(I)}\) ) into \(F_{m}\) such that \(g(i) = f(i)\) for all \(i \in I\) (I, § 7, nos. 1, 2 and 7); Proposition 7 then follows from no. 6, Proposition 6.

Remarks. (1) We shall later define an isomorphism of \(\operatorname{Libas}_{\mathbf{A}}(\mathbf{I})\) onto the tensor algebra of the free module \(\mathbf{A}^{(I)}\) ( \(\S 5\) , no. 5) and also an isomorphism of \(\operatorname{Libasc}_{\mathbf{A}}(\mathbf{I})\) onto the symmetric algebra of \(\mathbf{A}^{(I)}\) ( \(\S 6\) , no. 6).

\begin{enumerate}
\def\labelenumi{(\arabic{enumi})}
\setcounter{enumi}{1}
\item
  Let \(\rho\) be a unital homomorphism of A into a commutative ring B. As has been seen (§ 2, no. 6), an isomorphism \(\sigma\) is derived from \(\rho\) of \(\operatorname{Lib}_{\mathbf{B}}(\mathbf{I})\) (resp. \(\operatorname{Libas}_{\mathbf{B}}(\mathbf{I})\) , resp. \(\operatorname{Libasc}_{\mathbf{B}}(\mathbf{I})\) ) onto the algebra \((\operatorname{Lib}_{\mathbf{A}}(\mathbf{I}))_{(\mathbf{B})}\) (resp. \((\operatorname{Libas}_{\mathbf{A}}(\mathbf{I}))_{(\mathbf{B})}\) , resp. \((\operatorname{Libasc}_{\mathbf{A}}(\mathbf{I}))_{(\mathbf{B})}\) ) obtained by extending the scalars to B by means of \(\rho\) ; if \(\mathbf{X}_i^{\mathbf{A}}\) , \(\mathbf{X}_i^{\mathbf{B}}\) are the indeterminates of index \(i\) corresponding respectively to A and B, then \(\sigma(\mathbf{X}_i^{\mathbf{B}}) = \mathbf{X}_i^{\mathbf{A}} \otimes 1\) .
\item
  Let \(J\) be a subset of \(I\) ; we know that \(M(J)\) is identified with a stable subset of the magma \(M(I)\) and hence (no. 6) \(\operatorname{Lib}_{A}(J)\) is canonically identified with a subalgebra of \(\operatorname{Lib}_{A}(I)\) , generated by the \(X_i\) such that \(i \in J\) ; it is said that only the indeterminates of indices belonging to \(J\) occur in an element of \(\operatorname{Lib}_{A}(J)\) . The definition given in no. 6 of the algebra of a magma shows that \(\operatorname{Lib}_{A}(I)\) is the union of the directed family of subalgebras \(\operatorname{Lib}_{A}(J)\) when \(J\) runs through the set of finite subsets of \(I\) . There are analogous results for \(\operatorname{Libas}_{A}(I)\) and \(\operatorname{Libasc}_{A}(I)\) .
\item
  With element \(s\) of \(\mathbf{M}(\mathbf{I})\) (resp. \(\mathbf{Mo}(\mathbf{I})\) , resp. \(\mathbf{N}^{(1)}\) ) is associated its length \(l(s)\) , which is an integer \(\geqslant 1\) (resp. \(\geqslant 0\) ) such that \(l(ss') = l(s) + l(s')\) (I, §7, nos. 1, 2 and 7). If \(e_s\) is the element of \(\operatorname{Lib}_{\mathbf{A}}(\mathbf{I})\) (resp. \(\operatorname{Libas}_{\mathbf{A}}(\mathbf{I})\) , resp. \(\operatorname{Libasc}_{\mathbf{A}}(\mathbf{I})\) ) corresponding to \(s\) , the total degree (or simply the degree of an element \(x = \sum_{s} \alpha_s e_s \neq 0\) of \(\operatorname{Lib}_{\mathbf{A}}(\mathbf{I})\) (resp. \(\operatorname{Libas}_{\mathbf{A}}(\mathbf{A})\) , resp. \(\operatorname{Libasc}_{\mathbf{A}}(\mathbf{I})\) ) is the greatest of the numbers \(l(s)\) when \(s\) runs through the (non-empty by hypothesis) set of elements such that \(\alpha_s \neq 0\) . For example, if \(i, j, k\) are three distinct elements of \(\mathbf{I}\) , the element \((\mathrm{X}_t(\mathrm{X}_j\mathrm{X}_k))\mathrm{X}_t - (\mathrm{X}_t\mathrm{X}_j)(\mathrm{X}_k\mathrm{X}_l)\) is an element \(\neq 0\) of total degree 4 in \(\operatorname{Lib}_{\mathbf{A}}(\mathbf{I})\) .
\end{enumerate}

